\noindent\hspace*{2mm}%
\begin{minipage}[t]{\dimexpr\linewidth/2-4mm\relax}
\vspace{0pt}
\textbf{Versione quadratica}
\begin{unipdcode}[language=Python,
  basicstyle=\ttfamily\scriptsize,
  numbers=none,
  xleftmargin=0pt,
  xrightmargin=0pt,
  framexleftmargin=0.3em]
def somme_prefisse_quadratiche(valori):
    n = len(valori)              # 1
    risultato = [0] * n          # n + 1
    for i in range(n):           # n iterazioni
        somma = 0                # 1
        for j in range(i + 1):   # i + 1 iterazioni
            somma += valori[j]   # 2 2 (lettura, aggiornamento)
        risultato[i] = somma     # 1
    return risultato             # escluso
\end{unipdcode}
\end{minipage}%
\hspace*{4mm}%
\begin{minipage}[t]{\dimexpr\linewidth/2-4mm\relax}
\vspace{0pt}
\textbf{Versione lineare}
\begin{unipdcode}[language=Python,
  basicstyle=\ttfamily\scriptsize,
  numbers=none,
  xleftmargin=0pt,
  xrightmargin=0pt,
  framexleftmargin=0.3em]
def somme_prefisse_lineari(valori):
    n = len(valori)              # 1
    risultato = [0] * n          # n + 1
    somma = 0                    # 1
    for i in range(n):           # n iterazioni
        somma += valori[i]       # 2 (lettura, aggiornamento)
        risultato[i] = somma     # 1
    return risultato             # escluso
\end{unipdcode}
\end{minipage}

